\reading{MR-14}{The Art of Term Structure Models: Drift}{strike}{7}
% ==========================================================

\begin{mrkeybox}[title={Learning objectives in this reading}]
\small
\begin{enumerate}[label=\textbf{\alph*.},leftmargin=1.7em]
\item Construct and describe the effectiveness of a short-term interest rate tree assuming
      normally distributed rates, both with and without drift.
\item Calculate the short-term rate change and standard deviation of the rate change using a
      model with normally distributed rates and no drift.
\item Describe methods for addressing the possibility of negative short-term rates in term
      structure models.
\item Construct a short-term rate tree under the Ho-Lee Model with time-dependent drift.
\item Describe uses and benefits of the arbitrage-free models and assess the issue of fitting
      models to market prices.
\item Describe the process of constructing a simple and recombining tree for a short-term rate
      under the Vasicek Model with mean reversion.
\item Calculate the Vasicek Model rate change, standard deviation of the rate change, expected
      rate in $T$ years, and half-life.
\item Describe the effectiveness of the Vasicek Model.
\end{enumerate}
\end{mrkeybox}

\begin{mrdefbox}[title={Notation used throughout this reading}]
\begin{itemize}
\item $dr$ = change in interest rates over a small time interval $dt$
\item $dt$ = small time interval, measured in years (one month $= 1/12$)
\item $\sigma$ = annual basis point volatility of rate changes
\item $dw$ = normally distributed random variable with mean 0 and standard deviation
      $\sqrt{dt}$; for one month, $\sqrt{1/12} = 0.2887$
\item $\lambda$ = drift term added to the interest rate
\item $k$ = a parameter that measures the speed of reversion adjustment
\item $\theta$ = long-run value of the short-term rate assuming risk neutrality
\end{itemize}
\end{mrdefbox}

% ---------------------------------------------------------
\lo{MR-14 a}{Construct and describe the effectiveness of a short-term interest rate tree assuming normally distributed rates, both with and without drift.}

\subsection{Model 1 --- no drift}
Model 1 is assumed to have an interest rate term structure with \term{no drift}.

\begin{mrfmlbox}
\[ dr \;=\; \sigma\, dw
   \qquad\qquad
   \text{tree: } \begin{cases} r_0 + \sigma\sqrt{dt} \\[2pt] r_0 - \sigma\sqrt{dt} \end{cases} \]
\mrvk{with no drift term, the \textit{expected} change in the rate is zero. All that happens over
a step is a symmetric $\pm\sigma\sqrt{dt}$ shock.}{}
\end{mrfmlbox}

\begin{mrexbox}[title={Example --- calculating $dr$}]
Start with an initial rate of 6\%. Assume an annual volatility of 1.20\%. After one month, with
a random variable $dw$ of 0.2 drawn from a normal distribution:
\tcblower
\small
\[ dr \;=\; 1.20\% \times 0.2 \;=\; 0.24\% \;=\; \mathbf{24 \text{ basis points}} \]
The new rate after one month will be \textbf{6.24\%}.
\end{mrexbox}

% ---------------------------------------------------------
\lo{MR-14 b}{Calculate the short-term rate change and standard deviation of the rate change using a model with normally distributed rates and no drift.}

\begin{center}
\begin{tikzpicture}[every node/.style={font=\scriptsize\sffamily},
  bx/.style={rectangle, rounded corners=2pt, draw=FRMNavy!45, fill=PaleNavy,
    line width=0.6pt, inner sep=3.5pt, align=center},
  nx/.style={rectangle, rounded corners=2pt, draw=FRMGreen!50, fill=PaleGreen,
    line width=0.6pt, inner sep=3.5pt, minimum width=1.25cm, align=center}]
% symbolic
\begin{scope}
  \node[bx] (r0) at (0,0) {$r_0$};
  \node[bx] (u) at (2.0,1.1) {$r_0 + \sigma\sqrt{dt}$};
  \node[bx] (d) at (2.0,-1.1) {$r_0 - \sigma\sqrt{dt}$};
  \node[bx] (uu) at (4.8,2.0) {$r_0 + 2\sigma\sqrt{dt}$};
  \node[bx] (mm) at (4.8,0) {$r_0$};
  \node[bx] (dd) at (4.8,-2.0) {$r_0 - 2\sigma\sqrt{dt}$};
  \foreach \p/\q in {r0/u, r0/d, u/uu, u/mm, d/mm, d/dd}{
    \draw[FRMNavy!55, line width=0.7pt] (\p) -- node[sloped, above, text=black]{0.5} (\q);}
\end{scope}
% numeric
\begin{scope}[shift={(8.8,0)}]
  \node[nx] (r0) at (0,0) {6.000\%};
  \node[nx] (u) at (1.9,1.1) {6.346\%};
  \node[nx] (d) at (1.9,-1.1) {5.654\%};
  \node[nx] (uu) at (4.1,2.0) {6.692\%};
  \node[nx] (mm) at (4.1,0) {6.000\%};
  \node[nx] (dd) at (4.1,-2.0) {5.308\%};
  \foreach \p/\q in {r0/u, r0/d, u/uu, u/mm, d/mm, d/dd}{
    \draw[FRMGreen!55, line width=0.7pt] (\p) -- node[sloped, above, text=black]{0.5} (\q);}
\end{scope}
\end{tikzpicture}
\end{center}
\figcap{Model 1 recombines: an up-then-down move and a down-then-up move both land back on
$r_0$. The rate can drift arbitrarily far in either direction over many steps, including below
zero, which is the model's central weakness.}

\begin{mrexbox}[title={Example --- the standard deviation of the rate change}]
The current short-term rate is 6.18\%, volatility is 113 basis points per year, and the interval
is one month.
\tcblower
\small
\[ \sigma\sqrt{dt} \;=\; 1.13\% \times \sqrt{\tfrac{1}{12}} \;=\; 1.13\% \times 0.2887
   \;=\; \mathbf{0.326\%} \text{, or 32.6 basis points per month} \]
So the two date-1 nodes are $6.18\% + 0.326\% = \mathbf{6.506\%}$ and
$6.18\% - 0.326\% = \mathbf{5.854\%}$.
\end{mrexbox}

\begin{mrtrapbox}[title={Model 1 effectiveness}]
\begin{enumerate}
\item \term{Lacks flexibility}, but the tree is \term{recombining}.
\item \term{Term structure of interest rates:} downward sloping, and rates can go negative, due
      to the no-drift assumption.
\item \term{Flat volatility term structure:} it predicts a flat term structure of volatility, in
      contrast to the observed hump-shaped pattern.
\item \term{Single-factor model:} it relies solely on the short-term rate, ignoring drift and
      time-dependent volatility.
\item Produces \term{parallel} yield curve shifts.
\end{enumerate}
\end{mrtrapbox}

% ---------------------------------------------------------
\lo{MR-14 c}{Describe methods for addressing the possibility of negative short-term rates in term structure models.}

A problem with Gaussian models is that the short-term rate can become \term{negative}. Two
solutions:

\begin{center}
\small
\begin{tabularx}{\linewidth}{@{}p{4.0cm} >{\raggedright\arraybackslash}X >{\raggedright\arraybackslash}X@{}}
\arrayrulecolor{FRMRule}\toprule
\rowcolor{FRMNavy} \hdr{Method} & \hdr{How it works} & \hdr{The catch}\\
\textbf{1. Use non-negative distributions}
  & Use probability distributions that are always non-negative, such as lognormal or
    chi-squared. This ensures rates can never go negative.
  & May introduce other characteristics --- skewness, or volatilities --- that do not accurately
    reflect real-world conditions.\\
\rowcolor{FRMSoft}
\textbf{2. Force negative rates to zero}
  & Adjust the interest rate tree to ``force'' negative rates to be zero, constraining the
    distribution from going below zero.
  & Preferred in \textit{very low} interest rate environments, and it impacts a
    \textbf{narrower} range of rates.\\
\bottomrule
\end{tabularx}
\end{center}

% ---------------------------------------------------------
\lo{MR-14 d}{Construct a short-term rate tree under the Ho-Lee Model with time-dependent drift.}

\subsection{First, Model 2 --- constant drift}
This term structure model incorporates a \term{positive drift}, representing a risk premium for
longer time horizons.

\begin{mrfmlbox}
\[ dr \;=\; \lambda\, dt \;+\; \sigma\, dw \]
\mrvk{$\lambda$ = the positive constant drift, which captures two things at once: (1) a
\textbf{risk premium} and (2) \textbf{expected return}.}{}
\end{mrfmlbox}

\begin{center}
\begin{tikzpicture}[every node/.style={font=\scriptsize\sffamily},
  bx/.style={rectangle, rounded corners=2pt, draw=FRMNavy!45, fill=PaleNavy,
    line width=0.6pt, inner sep=3.5pt, align=center}]
  \node[bx] (r0) at (0,0) {$r_0$};
  \node[bx] (u) at (2.9,1.2) {$r_0 + \lambda dt + \sigma\sqrt{dt}$};
  \node[bx] (d) at (2.9,-1.2) {$r_0 + \lambda dt - \sigma\sqrt{dt}$};
  \node[bx] (uu) at (7.4,2.2) {$r_0 + 2\lambda dt + 2\sigma\sqrt{dt}$};
  \node[bx] (mm) at (7.4,0) {$r_0 + 2\lambda dt$};
  \node[bx] (dd) at (7.4,-2.2) {$r_0 + 2\lambda dt - 2\sigma\sqrt{dt}$};
  \foreach \p/\q in {r0/u, r0/d, u/uu, u/mm, d/mm, d/dd}{
    \draw[FRMNavy!55, line width=0.7pt] (\p) -- node[sloped, above, text=black]{0.5} (\q);}
\end{tikzpicture}
\end{center}
\figcap{The whole tree shifts up by $\lambda dt$ each step. The middle node at date 2 carries a
\textit{single} value, $r_0 + 2\lambda dt$, so this tree recombines.}

\begin{mrkeybox}[title={Model 2 effectiveness}]
\begin{enumerate}
\item \term{Term structure of interest rates:} due to the drift factor, the term structure will
      reflect an \term{upward-sloping} curve.
\item The tree \term{recombines} --- only one rate is produced at the middle node in year two.
\item Researchers typically choose parameters $r_0$ and $\lambda$ based on the calibration of
      observed rates, which enhances the model's fit to market data.
\end{enumerate}
\end{mrkeybox}

\begin{mrtrapbox}[title={Limitation of a constant drift}]
\begin{itemize}
\item The estimated value of the drift could be relatively \term{high}, particularly if seen as
      a risk premium only.
\item If the drift is viewed as a combination of the \textit{risk premium} plus \textit{expected
      rate changes}, the model implies rates are expected to keep increasing --- rates in year 10
      will be higher than in year 9. This view is more suitable for the \term{short term}, as it
      is harder to justify substantial increases in expected rates in the long run.
\end{itemize}
\end{mrtrapbox}

\subsection{The Ho-Lee model --- time-dependent drift}
The Ho-Lee model introduces greater flexibility by allowing the drift to vary over time. Each
time period can have a different drift, which does not have to follow a consistent upward
pattern and can even be negative.

\begin{mrfmlbox}
\[ dr \;=\; \lambda_t\, dt \;+\; \sigma\, dw \]
\mrvk{$\lambda_t$ = a \textit{different} drift for each period. The drift still captures risk
premium and expected return, but it is no longer forced to be the same number every step.}{}
\end{mrfmlbox}

\begin{center}
\begin{tikzpicture}[every node/.style={font=\scriptsize\sffamily},
  bx/.style={rectangle, rounded corners=2pt, draw=FRMNavy!45, fill=PaleNavy,
    line width=0.6pt, inner sep=3.5pt, align=center}]
  \node[bx] (r0) at (0,0) {$r_0$};
  \node[bx] (u) at (2.9,1.2) {$r_0 + \lambda_1 dt + \sigma\sqrt{dt}$};
  \node[bx] (d) at (2.9,-1.2) {$r_0 + \lambda_1 dt - \sigma\sqrt{dt}$};
  \node[bx] (uu) at (7.8,2.2) {$r_0 + (\lambda_1+\lambda_2) dt + 2\sigma\sqrt{dt}$};
  \node[bx] (mm) at (7.8,0) {$r_0 + (\lambda_1+\lambda_2) dt$};
  \node[bx] (dd) at (7.8,-2.2) {$r_0 + (\lambda_1+\lambda_2) dt - 2\sigma\sqrt{dt}$};
  \foreach \p/\q in {r0/u, r0/d, u/uu, u/mm, d/mm, d/dd}{
    \draw[FRMNavy!55, line width=0.7pt] (\p) -- node[sloped, above, text=black]{0.5} (\q);}
\end{tikzpicture}
\end{center}

\begin{mrkeybox}
\begin{enumerate}[label=\textbf{\alph*)}]
\item If $\lambda_1 = \lambda_2$, the Ho-Lee model \term{reduces to Model 2}.
\item $\lambda_1$ and $\lambda_2$ are estimated from observed market prices --- they are
      \term{calibrated}.
\item Each drift can take a \term{negative} value as well.
\end{enumerate}
\end{mrkeybox}

% ---------------------------------------------------------
\lo{MR-14 e}{Describe uses and benefits of the arbitrage-free models and assess the issue of fitting models to market prices.}

\begin{center}
\small
\begin{tabularx}{\linewidth}{@{}p{3.3cm} >{\raggedright\arraybackslash}X@{}}
\arrayrulecolor{FRMRule}\toprule
\rowcolor{FRMNavy} \hdr{1. Arbitrage-free models} & \hdr{model price $=$ market price}\\
\textbf{When used} & When the goal is to accurately price \term{illiquid or customized}
  securities. These models ensure the model price aligns with the market price.\\
\rowcolor{FRMSoft}
\textbf{Useful for} & Pricing securities like on-the-run Treasury bonds, and options on bonds.\\
\textbf{Drawbacks} & If the original model assumptions do not match real market conditions,
  pricing accuracy can suffer. These models assume underlying market prices are \textit{accurate},
  but they can be thrown off by temporary external shocks --- a sudden oversupply impacting market
  prices, for instance.\\
\midrule
\rowcolor{FRMNavy} \hdr{2. Equilibrium models} & \hdr{relative analysis}\\
\textbf{When used} & More suitable for \term{relative analysis}, where the goal is to compare the
  value of one security to another.\\
\rowcolor{FRMSoft}
\textbf{Assumption} & In relative analysis it is assumed that both securities are already
  properly priced, which makes equilibrium models more appropriate.\\
\bottomrule
\end{tabularx}
\end{center}

% ---------------------------------------------------------
\lo{MR-14 f}{Describe the process of constructing a simple and recombining tree for a short-term rate under the Vasicek Model with mean reversion.}

\begin{itemize}
\item The Vasicek model assumes a \term{mean-reverting} process for the short-term interest rate,
      which is a common assumption for the level of interest rates.
\item Mean reversion is a reasonable assumption but clearly breaks down in periods of extremely
      high inflation --- hyperinflation --- or similar structural breaks.
\end{itemize}

\begin{mrfmlbox}
\[ dr \;=\; k\,(\theta - r)\,dt \;+\; \sigma\, dw
   \qquad\qquad
   \theta \;\approx\; r_1 + \frac{\lambda}{k} \]
\mrvk{$k$ = speed of reversion adjustment --- \textbf{high $k$ means faster mean reversion};\quad
$\theta$ = long-run value of the short-term rate assuming risk neutrality;\quad $r$ = current
interest rate level;\quad $r_1$ = the long-run \textit{true} rate of interest;\quad $\lambda$ =
the annual drift.}{%
\par\smallskip\textbf{\sffamily Read the drift term:} $k(\theta - r)$ is positive when the rate
is below its long-run level and negative when it is above. That sign flip is the mean
reversion.}
\end{mrfmlbox}

\begin{mrexbox}[title={Example --- $dr$ under the Vasicek model}]
$k = 0.03$, annual standard deviation 150 basis points, a true long-term interest rate of 6\%, a
current interest rate of 6.2\%, and annual drift of 0.36\%.
\tcblower
\small
\[ \theta \;\approx\; 6\% + \frac{0.36\%}{0.03} \;=\; \mathbf{18\%} \]
\[ dr \;=\; k(\theta - r)\,dt \;=\; 0.03\,(18\% - 6.2\%)\left(\tfrac{1}{12}\right)
   \;=\; \mathbf{0.0295\%} \]
\end{mrexbox}

\begin{mrtrapbox}[title={Remember}]
For an interest rate tree using Vasicek, \term{add} volatility for up moves and \term{subtract}
volatility for down moves:
\[ dr = k(\theta - r)dt + \sigma\,dw \ \text{ for up moves}
   \qquad
   dr = k(\theta - r)dt - \sigma\,dw \ \text{ for down moves} \]
\end{mrtrapbox}

\subsection{Building the tree}

\begin{mrexbox}[title={Step 1 --- the two date-1 nodes, from $r_0 = 6.200\%$}]
\small
\begin{align*}
\text{up: }\ & 6.200\% + \frac{(0.03)(18\% - 6.200\%)}{12} + \frac{1.5\%}{\sqrt{12}}
  = \mathbf{6.663\%}\\[3pt]
\text{down: }\ & 6.200\% + \frac{(0.03)(18\% - 6.200\%)}{12} - \frac{1.5\%}{\sqrt{12}}
  = \mathbf{5.796\%}
\end{align*}
\end{mrexbox}

\begin{mrexbox}[title={Step 2 --- the date-2 nodes, recomputing the drift at each date-1 rate}]
\small
From the up node at 6.663\%:
\begin{align*}
& 6.663\% + \frac{(0.03)(18\% - 6.663\%)}{12} + \frac{1.5\%}{\sqrt{12}} = \mathbf{7.124\%}\\[3pt]
& 6.663\% + \frac{(0.03)(18\% - 6.663\%)}{12} - \frac{1.5\%}{\sqrt{12}} = \mathbf{6.258\%}
\end{align*}
From the down node at 5.796\%, the same two steps give \textbf{6.260\%} and \textbf{5.394\%}.
\end{mrexbox}

\begin{center}
\begin{tikzpicture}[every node/.style={font=\scriptsize\sffamily},
  nx/.style={rectangle, rounded corners=2pt, draw=FRMNavy!45, fill=PaleNavy,
    line width=0.6pt, inner sep=3.5pt, minimum width=1.35cm, align=center},
  sp/.style={rectangle, rounded corners=2pt, draw=FRMRed!55, fill=PaleRed,
    line width=0.7pt, inner sep=3.5pt, minimum width=1.35cm, align=center}]
  \node[nx] (r0) at (0,0) {6.200\%};
  \node[nx] (u) at (2.6,1.4) {6.663\%};
  \node[nx] (d) at (2.6,-1.4) {5.796\%};
  \node[nx] (uu) at (5.6,2.4) {7.124\%};
  \node[sp] (m1) at (5.6,0.38) {6.258\%};
  \node[sp] (m2) at (5.6,-0.38) {6.260\%};
  \node[nx] (dd) at (5.6,-2.4) {5.394\%};
  \foreach \p/\q in {r0/u, r0/d, u/uu, u/m1, d/m2, d/dd}{
    \draw[FRMNavy!55, line width=0.7pt] (\p) -- node[sloped, above, text=black]{0.5} (\q);}
  \draw[FRMRed, {Stealth[length=4pt]}-{Stealth[length=4pt]}, line width=0.8pt]
    (6.55,0.38) -- (6.55,-0.38);
  \node[anchor=west, text=FRMRed, align=left, text width=4.1cm] at (6.8,0)
    {two \textit{different} middle rates, so the tree does \textbf{not} recombine};
\end{tikzpicture}
\end{center}
\figcap{Mean reversion is what breaks recombination. The drift applied from the up node is
computed at 6.663\% and the drift from the down node at 5.796\%, so the two paths to the middle
arrive two-tenths of a basis point apart.}

\begin{mrkeybox}[title={Steps to make the tree recombining}]
\begin{enumerate}
\item Take an \term{average} of the two middle nodes:
      $(6.258\% + 6.260\%) / 2 = \mathbf{6.259\%}$.
\item Remove the assumption of 50\% up and 50\% down movements, generically replacing them with
      $(p,\ 1-p)$ and $(q,\ 1-q)$.
\end{enumerate}
\end{mrkeybox}

% ---------------------------------------------------------
\lo{MR-14 g}{Calculate the Vasicek Model rate change, standard deviation of the rate change, expected rate in $T$ years, and half-life.}

\begin{mrfmlbox}
\textbf{Expected rate in $T$ years}
\[ E[r_T] \;=\; \theta - (\theta - r_0)\,e^{-kT}
   \qquad\text{equivalently}\qquad
   E[r_T] \;=\; r_0\,e^{-kT} + \theta\bigl(1 - e^{-kT}\bigr) \]
\mrvk{$r_0$ = the current short-term rate;\quad $\theta$ = the long-run mean-reverting level;\quad
$k$ = the speed of mean reversion;\quad $T$ = the horizon in years. The gap between today's rate
and the long-run level decays exponentially at rate $k$.}{}

\medskip
\textbf{Half-life}
\[ \tau \;=\; \frac{\ln(2)}{k} \]
\mrvk{$\tau$ = the time it takes for the difference between the current rate and the long-term
mean to be reduced by \textit{half}.}{}
\end{mrfmlbox}

\begin{mrexbox}[title={Example --- expected rate in 10 years}]
Current short-term rate $r = 6.2\%$, mean reversion parameter $k = 0.03$, long-term
mean-reverting level $\theta = 18\%$.
\tcblower
\small
\textbf{1. Difference.} The gap between the current rate and the long-term average is
$18\% - 6.2\% = 11.8\%$.

\textbf{2. Decay.} That difference decays exponentially at the rate of mean reversion. After 10
years the remaining difference is
\[ 11.8\% \times e^{(-0.03 \times 10)} \;=\; \mathbf{8.74\%} \]

\textbf{3. Expected rate.} The long-term average \textit{minus} the decayed difference:
\[ 18\% - 8.74\% \;=\; \mathbf{9.26\%} \]
\end{mrexbox}

\begin{mrexbox}[title={Example continued --- half-life}]
\small
\[ \tau \;=\; \frac{\ln(2)}{0.03} \;=\; \frac{0.6931}{0.03} \;\approx\; \mathbf{23.1 \text{ years}} \]
It takes approximately 23.1 years for the gap between the current rate of 6.2\% and the
long-term average of 18\% to be halved --- that is, for the rate to reach around 12.1\%.
\end{mrexbox}

\begin{center}
\begin{tikzpicture}
\begin{axis}[
  width=11.0cm, height=5.4cm,
  xmin=0, xmax=60, ymin=5, ymax=19.5,
  xlabel={\scriptsize Years ahead, $T$}, ylabel={\scriptsize Expected short rate (\%)},
  xlabel style={font=\scriptsize}, ylabel style={font=\scriptsize},
  tick label style={font=\scriptsize},
  xtick={0,10,20,23.1,30,40,50,60},
  xticklabels={0,10,20,,30,40,50,60},
  ytick={6,9,12,15,18},
  axis line style={draw=black!55}, clip=false]
\addplot[FRMGold, dashed, line width=0.9pt, forget plot]
  coordinates {(0,18)(60,18)};
\node[anchor=east, font=\scriptsize\sffamily, text=FRMGold, fill=white, inner sep=1.5pt]
  at (axis cs:58,17.0) {$\theta = 18\%$ (long-run level)};
\addplot[FRMNavy, line width=1.3pt, smooth, domain=0:60, samples=120]
  {18 - 11.8*exp(-0.03*x)};
\draw[FRMRed, dashed, line width=0.7pt] (axis cs:23.1,5) -- (axis cs:23.1,12.1);
\draw[FRMRed, dashed, line width=0.7pt] (axis cs:0,12.1) -- (axis cs:23.1,12.1);
\node[circle, fill=FRMRed, inner sep=1.7pt] at (axis cs:23.1,12.1) {};
\node[anchor=west, font=\scriptsize\sffamily, text=FRMRed, fill=white, inner sep=1.5pt]
  at (axis cs:25,10.3) {half-life: $\tau = 23.1$ yrs $\Rightarrow$ rate $\approx 12.1\%$};
\node[circle, fill=FRMNavy, inner sep=1.6pt] at (axis cs:10,9.26) {};
\node[anchor=north west, font=\scriptsize\sffamily, text=FRMNavy, fill=white, inner sep=1.5pt]
  at (axis cs:11.5,8.7) {$T = 10 \Rightarrow 9.26\%$};
\node[anchor=west, font=\scriptsize\sffamily, text=black!60, fill=white, inner sep=1.5pt]
  at (axis cs:2.0,6.6) {$r_0 = 6.2\%$};
\end{axis}
\end{tikzpicture}
\end{center}
\figcap{The expected rate never reaches $\theta$; it approaches it asymptotically. The half-life
is where exactly half the original 11.8-point gap has closed, which is why the curve passes
through 12.1\% at 23.1 years. Both worked answers sit on this one curve.}

% ---------------------------------------------------------
\lo{MR-14 h}{Describe the effectiveness of the Vasicek Model.}

\begin{mrkeybox}[title={Vasicek model effectiveness}]
\begin{enumerate}[label=\textbf{\alph*)}]
\item \term{Term structure of volatility:} declining volatility is produced, where short-term
      volatility is relatively high while long-term volatility is comparatively low. Short-term
      volatility is overstated and long-term volatility understated.
\item \term{Term structure of interest rates:} specification is improved because the parameters
      $r_0$ and $\theta$ are calibrated to match observed market prices.
\item \term{Non-parallel shifts:} unlike some models, Vasicek can handle scenarios where interest
      rate shocks have varying impacts across different maturities. Short-term rates are affected
      more than long-term rates, and shocks dissipate as rates revert to the long-run mean.
      No-drift models, by contrast, produce parallel shifts affecting all rates equally.
\item \term{Effective for interest rate risk analysis} and yield curve modeling, where mean
      reversion is a critical factor.
\item \term{Interpretation of economic news:} a \textbf{larger} mean reversion parameter suggests
      quicker absorption of economic news --- short-term news means high $k$. A \textbf{smaller}
      parameter implies a more extended assimilation period --- longer-term news means low $k$.
\end{enumerate}
\end{mrkeybox}

\subsection{The three drift specifications side by side}

\begin{center}
\begin{tikzpicture}
\begin{axis}[
  width=11.0cm, height=6.1cm,
  xmin=0, xmax=30, ymin=4.5, ymax=13.5,
  xlabel={\scriptsize Years ahead}, ylabel={\scriptsize Expected short rate (\%)},
  xlabel style={font=\scriptsize}, ylabel style={font=\scriptsize},
  tick label style={font=\scriptsize},
  xtick={0,5,10,15,20,25,30}, ytick={5,7,9,11,13},
  legend style={font=\scriptsize\sffamily, draw=FRMRule, fill=white,
    at={(0.5,-0.30)}, anchor=north, legend columns=-1, column sep=10pt},
  axis line style={draw=black!55}, clip=false]
\addplot[black!60, line width=1.2pt, dashed, domain=0:30, samples=2] {6.2};
\addlegendentry{Model 1 --- no drift}
\addplot[FRMRed, line width=1.2pt, dotted, domain=0:30, samples=2] {6.2 + 0.22*x};
\addlegendentry{Model 2 --- constant drift $\lambda$}
\addplot[FRMNavy, line width=1.3pt, smooth, domain=0:30, samples=120]
  {18 - 11.8*exp(-0.03*x)};
\addlegendentry{Vasicek --- mean reverting to $\theta$}
\end{axis}
\end{tikzpicture}
\end{center}
\figcap{What the drift term actually does to the \textit{expected} path. Model 1 is flat forever.
Model 2 rises in a straight line and never stops, which is the limitation the notes flag --- it
is hard to justify rates rising without bound. Vasicek bends over toward $\theta$, which is why
it produces a declining volatility term structure rather than a flat one.}

\subsection{The four models in one table}

\begin{center}
\small
\begin{tabularx}{\linewidth}{@{}p{2.0cm} p{3.0cm} p{2.1cm} >{\raggedright\arraybackslash}X@{}}
\arrayrulecolor{FRMRule}\toprule
\rowcolor{FRMNavy} \hdr{Model} & \hdr{Process} & \hdr{Tree} & \hdr{Term structure produced}\\
\textbf{Model 1} & $dr = \sigma\,dw$ & Recombining
  & Downward sloping; flat volatility structure; parallel shifts; rates can go negative.\\
\rowcolor{FRMSoft}
\textbf{Model 2} & $dr = \lambda\,dt + \sigma\,dw$ & Recombining
  & Upward sloping, driven by the constant drift. Drift may be implausibly high over long
    horizons.\\
\textbf{Ho-Lee} & $dr = \lambda_t\,dt + \sigma\,dw$ & Recombining
  & Any shape --- each period's drift is calibrated separately and can be negative.
    Reduces to Model 2 if all $\lambda_t$ are equal.\\
\rowcolor{FRMSoft}
\textbf{Vasicek} & $dr = k(\theta - r)dt + \sigma\,dw$ & Non-recombining until the middle nodes
    are averaged
  & Declining volatility structure; non-parallel shifts; rates pulled toward $\theta$.\\
\bottomrule
\end{tabularx}
\end{center}
\figcap{The four models differ in one place only: the drift term. Model 1 has none, Model 2 has a
constant, Ho-Lee has a time-varying constant, and Vasicek makes the drift depend on where the
rate currently is. Everything else about the trees is identical.}


% ==========================================================
